% !TeX root = ../main.tex
% !TeX program = pdflatex
% !TeX encoding = UTF-8

\chapter{Bối cảnh bài toán}
\label{chap:problem-context}

\section{Mô tả tình huống giao thông được lựa chọn}

Đồ án nghiên cứu bài toán tìm đường giữa các địa điểm du lịch và đầu mối giao
thông tại thành phố Đà Lạt, tỉnh Lâm Đồng. Phạm vi dữ liệu gồm 25 địa điểm tiêu
biểu như Ga Đà Lạt, Chợ Đà Lạt, Hồ Xuân Hương, Thung Lũng Tình Yêu, Hồ Tuyền
Lâm và Bến xe Liên tỉnh Đà Lạt. Các địa điểm được chọn để tạo thành một mạng
lưới có cả khu trung tâm lẫn các điểm tham quan ngoại vi, qua đó thể hiện được
sự khác biệt về khoảng cách, loại đường, thời gian và điều kiện lưu thông.

Người dùng có thể chọn điểm xuất phát, điểm đến, kịch bản giao thông và tiêu
chí tối ưu. Ngoài bài toán một điểm đích, hệ thống còn hỗ trợ sắp xếp thứ tự
ghé nhiều địa điểm, phù hợp với một chuyến tham quan, giao hàng hoặc thực hiện
nhiều công việc trong cùng một hành trình.

\section{Bài toán thực tế cần giải quyết}

Địa hình Đà Lạt có nhiều dốc, đường quanh co và các tuyến dẫn đến điểm du lịch
nằm xa trung tâm. Lưu lượng giao thông còn thay đổi theo ngày thường, cuối
tuần, giờ cao điểm và điều kiện thời tiết. Vì vậy, tuyến ngắn nhất theo khoảng
cách chưa chắc là tuyến nhanh hoặc an toàn nhất. Đồ án tập trung giải quyết các
yêu cầu sau:

\begin{itemize}[itemsep=2pt]
  \item Tìm tuyến hợp lệ giữa hai địa điểm trên đồ thị có hướng.
  \item Tối ưu theo khoảng cách, thời gian hoặc chi phí tổng hợp có xét mức ùn
  tắc và rủi ro.
  \item Mô phỏng quá trình tìm kiếm để người học quan sát frontier, các đỉnh đã
  duyệt và tuyến cuối cùng.
  \item Sắp xếp thứ tự ghé nhiều địa điểm nhằm giảm tổng chi phí hành trình.
\end{itemize}

\section{Ý nghĩa của tối ưu hóa tuyến đường}

Tối ưu hóa tuyến đường giúp giảm thời gian di chuyển, nhiên liệu và chi phí vận
hành. Khi điều kiện đường thay đổi, hệ thống có thể đánh giá lại trọng số cạnh
để tránh các đoạn ùn tắc, mưa lớn, sương mù hoặc tạm đóng. Với nhiều điểm dừng,
việc tối ưu thứ tự ghé thăm còn tránh những lượt đi vòng không cần thiết.

Về mặt học thuật, bài toán là môi trường trực quan để so sánh các chiến lược
tìm kiếm không thông tin, có thông tin và tìm kiếm cục bộ. Mạng đường được mô
hình hóa thành đồ thị, còn yêu cầu tìm tuyến trở thành một bài toán tìm kiếm
trạng thái được trình bày trong chương tiếp theo.
